\section{Ist-Zustand}
\label{sec:ist-zustand}
\subsection{App-Steuerung}
Die App besitzt bereits eine theoretische Steuerung über die Bewegungssensoren des Smartphones. Dafür wird in gyroscope\_data.dart die Bibliothek motion\_sensors verwendet. Das Widget getGyroscope liest bereits die Werte für Pitch, Roll und Yaw aus.
\\
Sobald neue Sensordaten eintreffen, werden Pitch und Roll auf den Bereich von -1.0 bis 1.0 begrenzt. Anschließend werden diese Werte direkt als Steuerbefehle über den WebSocketManager an das Fahrzeug gesendet. Dabei soll der Roll-Wert die Geschwindigkeit steuern und der Pitch-Wert für den Lenkwinkel verwendet werden. Der Yaw-Wert wird für die Steuerung nicht verwendet. Zusätzlich zeigt die App die aktuellen Sensorwerte für Pitch, Roll und Yaw auf der Oberfläche an.
\\
In diesem Stand ist die Funktion zwar im Code angelegt, funktioniert praktisch aber noch nicht. Ein zentrales Problem liegt darin, dass die App von motion\_sensors.orientation abhängig ist. Das Paket motion\_sensors ist jedoch veraltet und liefert keine Werte. Außerdem werden die Sensorwerte direkt weitergegeben, ohne sie vorher zu filtern oder zu glätten. Dadurch könnten bereits kleinste Bewegungen oder Messungenauigkeiten die Steuerung beeinflussen.
\\
Weitere Probleme sind, dass es keine Deadzone für unbeabsichtigte kleine Bewegungen gibt und keine Möglichkeit vorgesehen ist, die Empfindlichkeit der Steuerung anzupassen. Auch eine einfache Korrektur falsch herum reagierender Achsen ist nicht vorbereitet.

\subsection{Funktionierend}
\label{sec:funktionierend}

\subsection{Probleme}
\label{sec:probleme}
% ==== Beginn Beitrag (Autorin: Sonia Sinaci) ====
Aus Sicht der Sensorik und der Systeminfrastruktur bestanden zum Zeitpunkt der Ist-Analyse vier wesentliche, im Projektverlauf dokumentierte Probleme.
\\
Erstens wurden die \gls{lidar}-Daten nicht richtig verarbeitet. Eine operative Nutzung realer \gls{lidar}-Hardware war nicht implementiert: Die vorhandenen \gls{lidar}-bezogenen Komponenten arbeiteten ausschließlich auf statischen, vorab aufgezeichneten Beispieldaten, und es fehlten Treiber- und Kommunikationsschnittstellen zur Anbindung des physischen Sensors. Im Fahrbetrieb lieferte der Sensor damit keinerlei verwertbare Daten.
\\
Zweitens funktionierten alle Sensoren bis auf einen nicht. Nach der Aktualisierung des Betriebssystems auf dem Raspberry~Pi brach die gesamte GPIO-Ansteuerung, wodurch neben Motorsteuerung und Lenkung auch die Sensorik ausfiel. Der Ultraschallsensor lieferte zusätzlich aufgrund eines Pinkonflikts mit dem SPI-Subsystem keine gültigen Messwerte. Zeitweise stand somit nur ein einziger funktionsfähiger Sensor zur Verfügung. Die Analyse und Behebung dieser Ausfälle sowie der daraus abgeleitete Wechsel auf den \gls{lidar} als primären Sensor werden in Abschnitt~\ref{sec:datenverarbeitung_sonia} beschrieben.
\\
Drittens beeinträchtigte der zu große Container die Anwendung. Das damalige Container-Image enthielt neben den Laufzeitabhängigkeiten auch vollständige Entwicklungswerkzeuge und grafische Anwendungen. Auf dem ressourcenbeschränkten Raspberry~Pi führte dies zu langen Bau- und Ladezeiten, und die Anwendung funktionierte dadurch nicht richtig (Lösung siehe Kapitel~\ref{sec:containerisierung_systemstart}).
\\
Viertens war der Systemstart nicht intuitiv. Das Starten des Systems erforderte eine Abfolge manueller Schritte im Container, die erst verstanden und eingerichtet werden musste. Ohne dieses Vorwissen war das System nicht in Betrieb zu nehmen. Zugleich war die Schrittfolge fehleranfällig und führte zu schwer reproduzierbaren Zuständen (Lösung siehe Abschnitt~\ref{sec:deterministischer_systemstart}).
% ==== Ende Beitrag ====
